
The machine learning community has expressed concern about epistemic lock-in: the possibility that reliance on a narrow set of benchmark datasets creates self-reinforcing cycles where early methodological choices become embedded in research practice. This concern appears well-founded---highly cited papers share benchmark datasets with their references at rates far exceeding chance, and concentration on standard evaluation corpora remains measurable across top venues. Yet despite these surface indicators of rigidity, we find no evidence that the field is trapped. Benchmark standards propagate through citation networks, but they do not create the temporal feedback loops that would constitute genuine lock-in.

This distinction matters because it separates legitimate concerns about evaluation methodology from unfounded fears about irrecoverable scientific stagnation. A field where benchmark use correlates with citation patterns differs fundamentally from one where early concentration causally determines future diversity. The former describes co-evolution between community standards and research practice; the latter describes a path-dependent trap.

\subsubsection{1.1 The Problem in Three Layers}

At the surface level, benchmark concentration is observable. Any researcher attending NeurIPS, ICML, or ICLR notices the same datasets appearing repeatedly: standard evaluation corpora for vision, language understanding, and other domains. This repetition raises questions about whether the field evaluates methods on sufficiently diverse tasks.

Digging deeper, citation analysis reveals that benchmark choices are not independent across papers. Prior work by Engdahl (2024) documents through ethnographic methods how benchmark standards emerge through community processes. High-impact papers share substantially more evaluation overlap with their references than citation structure alone would predict---a pattern consistent with benchmark standards propagating through intellectual influence.

Yet the deepest question remains: does this propagation create a feedback loop? If concentration in year t causally reduces diversity in year t+1, and this reduced diversity further increases concentration in year t+2, the field would face genuine lock-in requiring intervention. If instead concentration and citation patterns co-evolve without causal feedback---both responding to common factors like dataset availability, community growth, or methodological shifts---then benchmark standardization reflects coordination rather than confinement.

\subsubsection{1.2 Our Approach}

We analyze 12,600 papers from NeurIPS, ICML, and ICLR spanning 2018--2024, using Papers With Code data that provides benchmark usage information. We operationalize concentration through the Herfindahl-Hirschman Index (HHI), computing $\text{HHI} = \sum_d s_d^2$ where $s_d$ represents the share of papers using dataset d (operationalized via task labels as a proxy).

Our analysis proceeds through six sub-hypotheses testing the mechanism chain from concentration measurement through propagation to temporal dynamics:

\begin{enumerate}
\item \textbf{H-E1}: HHI can be reliably computed across all venue-years
\item \textbf{H-M1}: HHI correlates with simpler concentration metrics (top-5 share)
\item \textbf{H-M2}: Prior-year HHI predicts individual paper benchmark adoption
\item \textbf{H-M3}: Citation-linked papers share benchmarks at higher rates than random pairs
\item \textbf{H-M4}: Low diversity correlates with benchmark concentration
\item \textbf{H-M5}: Prior-year HHI predicts subsequent diversity decline (temporal lock-in test)
\end{enumerate}

\subsubsection{1.3 Contributions}

This paper makes three contributions:

\begin{enumerate}
\item We provide a quantitative characterization of dataset concentration across major venues, establishing that while concentration is measurable, it varies meaningfully across venues and years.
\end{enumerate}

\begin{enumerate}
\item We document the mechanism by which benchmark standards propagate: citation networks carry evaluation choices, creating correlation in benchmark use between papers and their references far exceeding chance overlap.
\end{enumerate}

\begin{enumerate}
\item We test whether this propagation creates self-reinforcing lock-in through panel regression and Granger causality analysis, finding that it does not.
\end{enumerate}

